\section{一次データソースの構造化}
\label{sec:llm}

\subsection{課題と設計要件}

一次データソースを直接収集しようとすると，構造が不安定であること（HTMLはスキーマを持たず，提供元の改修でセレクタが予告なく変わる）と，記述が多様で値が欠けていることの2つが問題になる．
LLMはこの両者に適するが，そのままデータ経路に置くと前稿で確立した検知結果の再現性を損なう．
そこで次の4要件を課す．
\begin{list}{}{%
  \setlength{\leftmargin}{1zw}\setlength{\itemindent}{-1zw}%
  \setlength{\labelwidth}{0pt}\setlength{\labelsep}{0pt}%
  \setlength{\topsep}{2pt}\setlength{\itemsep}{0pt}%
  \setlength{\parsep}{0pt}\setlength{\partopsep}{0pt}}
\item \textbf{R1（決定性）}同一入力から常に同一出力が得られる．
\item \textbf{R2（検証可能性）}出力を機械的に検証でき，通らないものはDBに入らない．
\item \textbf{R3（保守コストの局所化）}構造変更への追従が全面的な書き換えを要しない．
\item \textbf{R4（監査可能性）}生成元を追跡できる．
\end{list}
共通する設計はLLMの非決定性を日々のデータ経路から追い出すことであり，LLMに何を出力させるかで\textbf{P1}（データを一度だけ生成）・\textbf{P2}（抽出器のコードを生成）・\textbf{P3}（入力ごとに生成）の3つに整理した．この順はLLMが日々のデータ経路に残る度合いの順である．LLMは機械可読形式の代替ではなく，それが無いか不完全な場合の補完手段である．

\subsection{P1：欠損データの生成}
\label{sec:p1}

構造はあるが値が欠けている場合，LLMにデータを出力させ，凍結資産としてコミットする．実行は一度きりで日々のfetchでは呼ばれない（R1）．
Fortinet CVRFは2022年より前の \texttt{product\_statuses} が空で，1,119件のうち465件が該当する（値は\ref{sec:method3}節の影響を受けない版と同じく注記の自然文にしかない）．CSAFは2024年以降しか存在せず，この期間は一次データソースから検知できない．
そこでFortinet自身が書いた4つの記述——アドバイザリの2種の注記，CNAとして公開するCVEレコード，旧来の手書きデータセット——を突き合わせさせ，補完表を生成した\cite{vdu-pr892}．
検証は出典間の相互照合による（R2）．バージョン単位の照合，注記中のバージョン字句の網羅検査，NVDとの突合により計118件の不一致を検出し，根拠を記録しながら個別に裁定した（R4）．
バージョン制約を持たない19行は製品全体に一致してしまうため個別に監査し，未監査の行や未知の製品名は抽出を失敗させる——黙って落とすことも黙って広げることもしない．
結果として465件中422件が検知条件を得た．

\subsection{P2：抽出器の生成}
\label{sec:p2}

構造が安定していれば，fetchとextractのコードをLLMに書かせ，日々の収集は生成済みコードのみで決定的に動作させられる（R1）．検証はゴールデンセットへの回帰テストによる（R2）．
\ref{sec:eval}章で用いるApple Security Releases\cite{vdu-pr914}とMicrosoftの更新プログラム記事\cite{vdu-pr919}がこれにあたる．
このパターンの作業量は，構造の解析ではなく例外の発見と確定が占める．
Appleでは2003年以前のアーカイブ形式や公開が停止された記事の扱いが大半で，合計1,327行を要した．
このような網羅的な走査はLLMが得意とする作業であり，人手は生成された差分の承認に縮小される（R3）．
CIは取得件数の対前回比を閾値で監視してDBへの反映を止める．変更を再びLLMに与えて抽出器を直す運用は設計にとどまり，現状は人手で修正している．

\subsection{P3：HTMLからの直接変換}
\label{sec:p3}

構造が弱く鍵も欠けている場合，新規入力ごとにLLMにデータを出力させる．ここだけはデータ経路にLLMが残るため，変換結果をGitで持つことで決定性を作り（R1），型と検証器で検査する（R2）．
本稿が調べた範囲で最も極端なのはOpenSSHであった．
第一に，項目の境界が位置にしか現れない．セキュリティ情報のページ\cite{openssh-sec}の55件のエントリは \texttt{<li>} の中で日付・影響範囲・表題・本文が \texttt{<br>} で区切られているだけで \texttt{class} 属性も見出しも持たず，\texttt{<br>} の数は2個から7個までばらつくため，何番目の断片が影響範囲かをセレクタで固定できない．
第二に，鍵となるCVE IDがほとんど振られていない．本文にCVE IDを含むエントリは6件にとどまる．OpenSSHはCNAではなく，我々が観測した134件のCVEはMITREやRed Hatが採番している．アドバイザリを書く主体とCVEレコードを書く主体が別であることが欠落の理由である．
その帰結が\ref{sec:c1c3}節のCVE-2024-6387（regreSSHion）である．ページはこれを \texttt{8.5p1 to 9.7p1 (inclusive)} と正確な範囲付きで記述しながらCVE IDを書いていない．
検知条件は一次データソースだけで組み立てられるが，それが何と呼ばれている脆弱性なのかは分からない．

55件のエントリをすべて \texttt{vuls-data-raw} 形式へ変換した\cite{vdu-pr928}．
検査は2段ある．extractのコードは各バージョン文字列がOpenSSHのリリース表記かと範囲の上下限が高々1つで空でないかを検べ，通らなければ抽出を失敗させる．一方，出力の版とCVE IDが原文に逐語で現れるか，変換元の \texttt{<li>} を連結すると元のリストが復元されるかは，手順書がモデルに課す自己検査であり機械的に強制していない．原文に無い値の混入を止める検査こそがR2を満たしていない．
CVE IDの紐付けは一次データソースだけでは6件（10.9\%）にとどまるが，MITREのCVEレコードを参照することで45件（81.8\%）まで上がった．
一次データソースは影響範囲を持ちながら索引を持たず，MITREは索引を持ちながら影響範囲の精度を欠くという相補性が，実測として確認された．
ただしこの紐付けは検知件数を1件も変えていない．
生成元の追跡（R4）は，fetchがページを \texttt{origin/} に逐語で保存し，記録ごとに変換元の \texttt{<li>} を持たせることで満たす．変換の陳腐化——ページが変わったのに \texttt{raw/} を作り直していない状態——はCIで検出する\cite{vdb-pr224}．決定性（R1）は，変換結果である \texttt{raw/} をGitで持ち，fetchがそれを保存することで満たす．ページが変わらない限りモデルは走らず，出力は取得し直しても変わらない．残る不確かさは，ページが変わったときの再変換が同じ出力になる保証がないことである．
